\vspace{14pt}
\begin{center}
{\bfseries Period-Eight Spectral Structure and Counterexamples for Signed Circulant Graphs\par}
\vspace{6pt}
\ifdefined\TargetAAnonymous Anonymous\else [AUTHOR NAME]\fi
\end{center}

\noindent\textbf{Abstract:}
For the signed circulant graph \(C_n(1,2)\), a recent conjecture asserts that the
alternating-triangle-flux class with negative Hamilton-cycle holonomy minimizes
the spectral radius for every even \(n\geq 8\). We disprove this conjecture and
determine its smallest failure. A complete exact search verifies the proposed
bound for every even \(8\leq n\leq 30\), whereas an explicit signing at \(n=32\)
has strictly smaller spectral radius. The finite minimality statement is
computer-assisted: its canonical representative sets are checked by two
independent enumeration routes, and every spectral decision uses exact integer
or rational arithmetic. The counterexample arises from a period-eight flux
word and extends to every \(n\geq32\) divisible by eight. Floquet reduction gives
an exact quartic for the squared bands and the sharp infinite-volume edge
\(\eta=4+\sqrt{10+2\sqrt{5}}\). The phase has a local structural explanation:
after the signed operator is squared, negative quadrilateral flux cancels the
associated odd-displacement couplings, while positive flux activates localized
defects. Closed-walk moments then yield a complete period-eight trichotomy.
Antipodal defect pairs are exactly the sub-eight phases, the all-negative phase
has squared radius eight, and every other legal phase lies above eight. For
arbitrary period, the first three exact moments give necessary defect-density
and clustering inequalities. Finally, an exact classification proves that the
period-eight phase is the unique minimizer, modulo the natural operator
equivalences, among phases of primitive Hamilton-gauge period at most sixteen.\par

\noindent\textbf{Keywords:} signed circulant graph; spectral radius; switching
equivalence; flux phase; Floquet theory; closed-walk moments; computer-assisted proof
